% Chinese (中文) -- what the reading editions' preamble leaves to the language.
% Input by lib/frank-preamble.tex as \input{\FrankLib/lang/\BookLang.tex};
% docs/languages.md section 6 lists the hooks every file here must define and
% section 12 the ones it may, and the manual's "Adding a language" chapter
% says in one line each what happens when one is missing.
%
% Simplified characters: everything zh teaches is written in them
% (docs/lang/zh.md), and this file names one font accordingly.
%
% THE CJK MODEL, MINUS THE READING.  This is lib/lang/ja.tex without the
% furigana: no word separator, so the seam between two chunks is glue and not
% a space, and a pass that is a DIRECTION rather than a font -- the text set
% vertically.  What Japanese has and Chinese has not is a second layer of
% sound: kana is a reading written in the language's own script, and Chinese
% has no such script, so pinyin is the transliteration and goes where every
% other language's does, in the `tr' field -- \ch's third argument, and the
% reason \FrankHasReading is false here and the gloss block carries three
% lines and not four.  Nothing comes off the text either (the registry's `strip` is
% null and Chinese writes no marks over its characters), and while pass 1 was
% the text alone that meant no bare pass: it would have repeated pass 1.  A
% chunk that carries its words, \chw, sets pinyin over each of them in pass 1
% -- the per-word ruby Japanese sets in kana, \FrankRuby below -- so pass 1
% has help to take off again, and a passage comes FIVE times: the sentence
% with pinyin over its words, the pinyin alone, the chunks, the plain text,
% the vertical setting.  (A book still written in \ch has no pinyin in pass 1,
% and its plain pass repeats it until its chunks are given their words.)
%
% As in ja.tex, everything below is done with babel and fontspec alone --
% there is no luatexja on this machine and the editions are built with
% LuaLaTeX: babel's `chinese' locale with intraspace glue for the line breaks
% (Chinese has no spaces to break at), and the vertical pass as a rotated box
% in a font loaded with the OpenType `vert' feature.

% ---- the babel language and the switches ----------------------------------
% \FrankLangName is what \foreignlanguage and \begin{otherlanguage} are given
% everywhere in the core preamble: babel must know this name, or every passage
% comes out in the roman with a "Unknown language" error.
\newcommand{\FrankLangName}{chinese}
\FrankRTLfalse           % chunk left, gloss right; \pw needs no bidi isolate
\FrankHasBaretrue        % pass 4 is the text without the pinyin pass 1 sets over its words
\FrankHasAltfalse        % no second face: the last pass is the vertical one
\FrankHasReadingfalse    % pinyin is the transliteration, not a reading line
\FrankHasVerttrue        % so pass 5 is \FrankVertPass, not the font pass
\FrankHasAloudtrue       % pass 2 is the pinyin alone, chunk after chunk

% ---- fonts ------------------------------------------------------------------
% No face is NAMED here: \FrankPickFont walks the registry's tex.main and then
% tex.main_fallbacks and takes the first one installed, so the .tex and the
% .json cannot drift apart (docs/languages.md section 12).  A bundled face --
% one that travels in lib/fonts/ -- is found because build.sh puts that
% directory on OSFONTDIR; nothing CJK travels with the toolbox, the files are
% too large, so this one is the system's Noto Serif CJK SC.  SC is the
% simplified cut, and it is the only thing in this file that would change for
% a traditional edition.
% intraspace=0 .1 0 is the glue babel puts between two CJK characters: no
% natural width, a tenth of an em of stretch, no shrink -- what lets a line
% break anywhere and still be justified.  Renderer=HarfBuzz as for the Arabic
% script and for Japanese, so one engine shapes every language.  Script=CJK is
% fontspec's name for the hani script.
% Language is fontspec's OpenType language tag, and fontspec knows no plain
% `Chinese': the four it declares are Chinese Simplified (ZHS), Chinese
% Traditional, Chinese Hong Kong and Chinese Phonetic.  The registry's
% tex.language carries the name in full for that reason -- the row said
% `Chinese' when it was scaffolded, and every book died at the preamble with
% "Key 'fontspec-opentype/Language' accepts only a fixed set of choices".
\babelprovide[import=zh,onchar=ids fonts,intraspace=0 .1 0]{chinese}
\FrankPickFont{\FrankMainFontName}{zh}{main}{main_fallbacks}
\babelfont[chinese]{rm}[Script=CJK,Language=Chinese Simplified,Renderer=HarfBuzz]{\FrankMainFontName}
% The vertical face: the main font again, with the `vert' and `vrt2' features
% on, so that the glyphs which change shape in vertical text -- the brackets,
% the comma, the full stop, the ellipsis -- are the vertical forms.  Set
% horizontally in this font and turned clockwise, each glyph then stands
% upright in its column.  This is ja.tex's arrangement unchanged, and for the
% same reason: HarfBuzz honours the features and leaves the text layer alone,
% so pdftotext gives back the characters the book was written with.
\newfontfamily\FrankVertFont[Script=CJK,Language=Chinese Simplified,Renderer=HarfBuzz,%
  RawFeature={+vert,+vrt2}]{\FrankMainFontName}
\usepackage{graphicx}   % \rotatebox and \resizebox for the vertical pass

% ---- the two Lua filters ---------------------------------------------------
% frank_strip(s): the bare form of a chunk.  Nothing is written over a Chinese
% character and nothing comes off, so this is the identity: the plain pass
% and the vertical pass set the text as it is written.
% frank_latin(s): the label's digits as ASCII.  The registry's `digits' for
% Chinese are the Latin ones, which is what a modern Chinese book numbers its
% chapters with, so this is the identity too.  Both are called for every
% language, so both must exist even when they do nothing.  No #, % or literal
% backslash in a \directlua body: TeX reads it first (NOTES section 9 of the
% Persian book).
\directlua{
  function frank_strip(s) return s end
  function frank_latin(s) return s end
}

% ---- the seam between two chunks --------------------------------------------
% The core joins the chunks of a sentence with \FrankChunkSep when it rebuilds
% pass 1 and the plain and vertical passes from the buffers.  A space, the default, is right for a
% language written with spaces; Chinese is not, and a space there would show
% as a gap the source does not have.  So the seam is glue of no width with the
% same stretch babel puts between two characters: invisible, and a place to
% break the line.
\newcommand{\FrankChunkSep}{\hskip 0pt plus .1em\relax}

% ---- the vocabulary line -----------------------------------------------------
% \vb{verb}{pinyin}{A了B}{pinyin}{A不B}{pinyin}{meaning}: A CHINESE VERB HAS NO
% FORMS.  Nothing is added to it and nothing comes off it: what a European
% verb says with an ending, Chinese says with a separate word standing next
% to the verb, and 吃 is 吃 in every person, number and time there is.  So
% almost every verb -- and every single character -- is glossed with \dw like
% any other headword, and the three-form model has nothing to put in its
% slots: its old labels, "verb" and "with 了", would have printed the headword
% again and a V了 that reads as a past tense, which zh.md forbids.
%
% What a reader DOES need beside two kinds of verb, and cannot see in the
% characters, is where it comes apart, and so \vb is for those two only:
%   split  a separable verb-object verb (离合词): the verb with 了 between its
%          halves, pinyin as zh.md writes a particle, joined to what precedes
%          it -- \vb{睡觉}{shuìjiào}{睡了觉}{shuìle jiào}{}{}{to sleep}.  Without
%          it 睡了一觉, 帮我的忙 and 见过面 cannot be traced back to a word at all.
%   can't  a verb with a resultative or directional complement: the potential
%          with 不 inside, one word in pinyin -- \vb{看见}{kànjiàn}{}{}{看不见}
%          {kànbujiàn}{to see}; 得 in the same place is "can", so it is not given.
% A verb is one or the other, never both, which is why \vb prints a pair only
% when its form is given (lib/frank-preamble.tex): the split verb leaves the
% third pair blank and the complement verb the second, and neither prints a
% label standing over nothing.  Chinese has no extras in the parenthesis.
% The editor's sidebar proposes these from the dictionary (lib/verbs/zh.py):
% the kind from Wiktionary's zh-verb type (vo, vc), where it parts from the
% entry's canonical row (睡⫽覺), and lib/lang/zh.verbs.json for the common
% verbs Wiktionary leaves unmarked (散步, 听懂, 回来) or marks wrongly (想到).
% The pair below is the registry's ("vb_labels": ["split", "can't"], with the
% three forms in words as "vb_forms"), and lib/newlang.py --check compares it
% with these two lines, so the two files must agree.
\newcommand{\FrankVbPres}{split}
\newcommand{\FrankVbPast}{can't}

% ---- pinyin over a word --------------------------------------------------------
% \FrankRuby{word}{pinyin}: what pass 1 and the chunk column set over each word
% of a \chw -- lib/lang/ja.tex's \jruby, with pinyin for kana.  The pinyin is
% small and upright, through \FrankRoman as every transliteration in the book
% is, centred over the word, and the pair is one unbreakable box, which is why
% the core puts \FrankWordSep between two words.  Pinyin is wider than the
% characters it reads (zhōngguó over 中国), so the box is as wide as the wider
% of the two -- and the pinyin is measured with a fifth of an em either side:
% a word's pinyin has spaces of its own (yí ge), and without that margin the
% pinyin of two words stood no further apart than two of its syllables.
% Measured first, as \jruby is: a pair wider than the measure -- the chunk
% column is a third of the line -- is set stacked instead, the pinyin on a
% line of its own and the word under it.
\newcommand{\FrankRuby}[2]{%
  \begingroup
    \setbox0\hbox{{\SzTrans\kern.2em\FrankRoman{#2}\kern.2em}}%
    \setbox2\hbox{#1}%
    \dimen0=\wd0 \ifdim\wd2>\dimen0 \dimen0=\wd2 \fi
    \ifdim\dimen0>\hsize
      \par\noindent{\SzTrans\FrankRoman{#2}\par}\nobreak\noindent #1%
    \else
      \leavevmode
      \vbox{\hbox to\dimen0{\hss\box0\hss}%
            \nointerlineskip\kern0.5pt
            \hbox to\dimen0{\hss\box2\hss}}%
    \fi
  \endgroup}
% \FrankAloudText{pinyin}: pass 2 sets each chunk's pinyin through \FrankRoman
% too; the text's closing punctuation after it stays in the Chinese face.
\newcommand{\FrankAloudText}[1]{\FrankRoman{#1}}

% ---- the vertical pass -------------------------------------------------------
% The text (\FrankBuf, the same buffer the plain pass sets) as
% 直排: set in a \parbox whose WIDTH is the column height, in the
% vertical-forms font, then turned clockwise with \rotatebox.  Line 1 becomes
% the rightmost column and the columns run right to left, which is the order a
% vertically set Chinese page is read in.
%
% HOW TALL THE COLUMNS ARE.  The turned box cannot be broken across pages, so
% its height decides how the pages fill.  Columns of a fixed height whatever
% the sentence's length follow the pass before them on the same page almost
% never, and leave the page before mostly white; so the sentence is measured
% first -- its natural width in the vertical face, \wd1, is its length as a
% single column -- and spread over as many columns as it needs at
% \FrankVertColumn, each column getting an equal share plus one em of slack so
% that the whole-character line breaks can always fill the last column from
% the others.  A short sentence is then a short box that sits under the
% plain pass; a long one has many columns of a comfortable height rather than few
% of a great one.  When that would want more columns than the measure holds (a
% column is one leading wide once turned) the columns grow towards
% \FrankVertHeight instead, and only a sentence too long even for that is
% scaled down to \linewidth.
% The column count is worked out in Lua because it is a ceiling and a floor of
% quotients of dimensions, which TeX's own arithmetic (\numexpr rounds) makes
% needlessly awkward: frank_vert_columns(width of the sentence, usable column
% height, usable tallest height, measure, leading) in sp.
% \raise removes the depth the rotation gives the box, so it sits on the page
% like a paragraph and the page breaker sees its whole height.
% Every number here was tuned on lib/lang/ja.tex's fixture edition and is kept
% unchanged: the two languages set the same square characters at the same size
% down the same column, and a value that balanced a Japanese page balances a
% Chinese one.
\newcommand{\FrankVertHeight}{0.78\textheight}   % the tallest a column may ever be
\newcommand{\FrankVertColumn}{0.5\textheight}    % the height the columns are balanced to
\directlua{
  function frank_vert_columns(w, column, tallest, measure, leading)
    local most = math.floor(measure / leading)
    local n = math.ceil(w / column)
    if n > most then n = math.ceil(w / tallest) end
    if n < 1 then n = 1 end
    return n
  end
}
\newcommand{\FrankVertPass}{%
  \par\addvspace{1.6ex}%
  \begingroup
    \FrankVertFont\SzPlain
    \setbox1\hbox{\foreignlanguage{\FrankLangName}{\FrankBuf}}%
    % the usable heights are the caps less the em of slack every column gets,
    % so a balanced column never exceeds its cap; \fontdimen6 is the em
    \dimen1=\dimexpr\FrankVertColumn-\fontdimen6\font\relax
    \dimen2=\dimexpr\FrankVertHeight-\fontdimen6\font\relax
    \count1=\directlua{tex.print(frank_vert_columns(\number\wd1, \number\dimen1,
      \number\dimen2, \number\linewidth, \number\baselineskip))}\relax
    \dimen3=\dimexpr\wd1/\count1+\fontdimen6\font\relax
    \setbox0\hbox{\rotatebox[origin=c]{-90}{%
      \parbox{\dimen3}{\FrankVertFont\SzPlain\color{ink}%
        \begin{otherlanguage}{\FrankLangName}\FrankBuf\par\end{otherlanguage}}}}%
    \ifdim\wd0>\linewidth
      \setbox0\hbox{\resizebox{\linewidth}{!}{\box0}}%
    \fi
    \setbox0\hbox{\raise\dp0\box0}%
    \noindent\hfill\box0\par
  \endgroup}

% ---- the "How to read this" page -----------------------------------------------
% Printed by lib/frank-frontmatter.tex, before the first chapter: the only
% place the reader is told what the passes are for.  \pw{...} sets one word of
% the language inside the English; it is defined by the core preamble, after
% this file, and only expanded when the page is set.
\newcommand{\FrankHowTo}{%
Every passage comes five times.

\smallskip
\textbf{First} the whole sentence as Chinese is written, character after
character with nothing between them, and its pinyin small above each word.
Try to read it. Get as far as you can and do not worry about what you cannot
get — the point is the attempt, made before any help arrives. The pinyin
stands over words and not over characters, because a Chinese word is usually
more than one character: \pw{中国} is one word written with two.

\smallskip
\textbf{Second} only the sound of it: the pinyin of each phrase, a gap between
one phrase and the next, and no characters at all. Read it aloud, tones and
all. The gaps are where the sentence divides.

\smallskip
\textbf{Third} the same sentence in those phrases, Chinese on the left, and on
the right three lines: how it sounds in pinyin, then how to look each word up,
then what it means. Now you find out how much you had right.

\smallskip
\textbf{Fourth} the same passage plain, with no pinyin anywhere — the way
Chinese is printed. You have read it three times, so it should carry you.

\smallskip
\textbf{Fifth} the same again set vertically, top to bottom, the columns
running from right to left — the way Chinese was set for two thousand years
and is still set on a title page, a scroll and a shop sign. You already know
what it says, which is the only comfortable way to learn a new direction.

\smallskip
The pinyin is the whole of the sound: there is no second spelling underneath
it, as there is in a Japanese edition, because Chinese writes no alphabet of
its own. The mark over a vowel is the tone, and the tone is part of the word —
\textit{m\=a} mother, \textit{m\'a} hemp, \textit{m\v{a}} horse,
\textit{m\`a} to scold. A syllable with no mark is said light and short
(\textit{de}, \textit{le}, \textit{ma}). \textit{q} is roughly the
\textit{ch} of \textit{cheese} and \textit{x} the \textit{sh} of \textit{she},
both said far forward; \textit{zh ch sh r} are the same sounds with the tongue
curled back; \textit{c} is \textit{ts} and \textit{z} is \textit{dz}.
\textit{\"u} is the vowel of French \textit{tu}, and after \textit{j q x y} it
is written \textit{u} because nothing else can stand there. An apostrophe only
divides syllables that would otherwise run together: \pw{西安} is
\textit{X\=i'\=an}, two syllables, not \textit{xian}.

\smallskip
Verbs are given as one form, because a Chinese verb has only one: nothing is
added to it for person, number or time. What a tense would say is said by the
words around the verb — \pw{了} that the thing was finished, \pw{过} that it
has been done at some time, \pw{着} that it is going on — and the meaning line
on the right is where you will find it. Two kinds of verb come apart in the
sentence, and for those two the vocabulary shows how: \textit{split} is a verb
made of a verb and its object, which takes \pw{了} between its halves
(\pw{睡觉}, \pw{睡了觉}); \textit{can't} is a verb with a result or a
direction after it, which takes \pw{不} between them (\pw{看见}, \pw{看不见},
cannot see) — and \pw{得} in that place means \textit{can}. Nouns are
counted with a measure word standing between the number and the noun, and the
vocabulary names it whenever it is not the general \pw{个}: \pw{一个人} is
\textit{one}, then the measure word, then \textit{person}.
}
